\documentclass[11pt]{article}
\usepackage[a4paper,margin=21mm]{geometry}
\usepackage{fontspec}
\setmainfont{Latin Modern Roman}
\newfontfamily\CJKfont{Noto Serif CJK SC}
\XeTeXlinebreaklocale "zh"
\XeTeXlinebreakskip=0pt plus 1pt
\usepackage{amsmath,amssymb,amsthm,amscd,graphicx,color}
\usepackage{xurl}
\usepackage[unicode,hidelinks]{hyperref}
\allowdisplaybreaks
\setlength\emergencystretch{3em}
\numberwithin{equation}{section}
\newtheorem{Theorem}{Theorem}[section]
\newtheorem{Corollary}[Theorem]{Corollary}
\newtheorem{Lemma}[Theorem]{Lemma}
\newtheorem{Proposition}[Theorem]{Proposition}
 { \theoremstyle{definition}
\newtheorem{Definition}[Theorem]{Definition}
\newtheorem{Note}[Theorem]{Note}
\newtheorem{Example}[Theorem]{Example}
\newtheorem{Remark}[Theorem]{Remark} }

\numberwithin{equation}{section}

\usepackage[mathscr]{eucal}
\newcommand{\ZZ}{\mathbb{Z}}

\makeatletter\newlabel{CLSXYthm}{{1.1}{}{Original source locator}{}{}}\makeatother
\begin{document}
\begin{center}\Large Uniform shifted-binomial finite Bailey-pair transformation\end{center}
\noindent\textbf{Stable ID:} 556\quad\textbf{Author:} Shashank Kanade; Jeremy Lovejoy\par
\noindent\textbf{Work:} Bailey Pairs and an Identity of Chern–Li–Stanton–Xue–Yee\par
\noindent\textbf{Source:} \url{https://sigma-journal.com/2025/021/}\par
\noindent\textbf{Version:} Official SIGMA 21 (2025), article 021, published 2025-03-29; journal PDF and publisher native TeX archive acquired 2026-09-30; exact original bytes pinned by SHA256\par
\noindent\textbf{Rights:} CC BY 4.0. Authors retain copyright. \url{https://creativecommons.org/licenses/by/4.0/}. Reuse and typesetting adaptation; original language and mathematical content preserved.\par
\noindent\textbf{Difficulty (prior selection preserved):} {\CJKfont H2: The research-level core is the uniform construction for arbitrary m and shift position a: compose Bailey iterations with a raising transform q→q² and a lattice lowering q²→q so that an adjacent f-term cancels while exactly one internal q-binomial acquires the delta shift. The nonlocal multisum indexing, the n=0 lattice boundary, and separate a=0/a=m endpoints must be coordinated. This is not a definition check, a single named-theorem invocation, or ordinary qualifying-exam summation; all six construction stages are retained.}\par
\medskip\noindent\textbf{Editorial boundary note (not author text):} Complete original Lemma 3.1 as the uniform shifted-binomial finite Bailey-pair transformation. Retains full q-product/binomial and Bailey-pair conventions, the three author-cited Bailey lemma/lattice inputs, all six transformation steps, boundary relation at n=0, both a=m and a=0 cases, and final coefficient rewrite. This one existing unit is the general transformation itself; later applications/identity specializations are not separate counts. The author positions it as the key Bailey lemma for the paper’s main theorem; it has its own arbitrary-input Bailey-pair statement and is not merely one specialized scalar calculation.\par
\noindent\textbf{External original reference CLSXYthm:} Original cited Chern–Li–Stanton–Xue–Yee identity, source lines 93–97; mentioned only in the background discussion about which Bailey lemma is used later.\par
\medskip\hrule\medskip\noindent\textbf{Author text begins below.}\par
\setcounter{section}{1}\subsection*{Author q-series notation}
% Original source lines 75--90
Recall the usual $q$-series notation,
\begin{equation*}
	(a_1,a_2,\dots ,a_k)_{\infty} = (a_1,a_2,\dots ,a_k;q)_{\infty} = \prod_{j=0}^{\infty} \bigl(1-a_1q^j\bigr)\bigl(1-a_2q^j\bigr) \cdots \bigl(1-a_kq^j\bigr)
\end{equation*}
and
\begin{equation*}
	(a_1,a_2,\dots ,a_k)_n = (a_1,a_2,\dots ,a_k;q)_n = \prod_{j=0}^{n-1} \bigl(1-a_1q^j\bigr)\bigl(1-a_2q^j\bigr) \cdots \bigl(1-a_kq^j\bigr),
\end{equation*}
valid for $n \geq 0$, along with the $q$-binomial coefficient
\begin{equation} \label{eqn:qbinom}
	\begin{bmatrix} n \\ k \end{bmatrix} = \begin{bmatrix} n \\ k \end{bmatrix}_q =
	\begin{cases}
		\displaystyle\frac{(q)_n}{(q)_{n-k}(q)_k} & \text{if $0 \leq k \leq n$}, \\
		0 & \text{otherwise}.
	\end{cases}
\end{equation}

\setcounter{section}{2}\setcounter{Theorem}{0}\setcounter{equation}{0}\subsection*{Author Bailey-pair prerequisites}
% Original source lines 241--304
In this section, we review the necessary background on Bailey pairs. A Bailey pair relative to~$x$~\cite{Andmultiple} is a pair of sequences $(\alpha_n,\beta_n)$,
$n\in \ZZ_{\geq 0}$,
satisfying
\begin{equation*} %\label{pairdef}
	\beta_n = \sum_{k=0}^n \frac{\alpha_k}{(q)_{n-k}(xq)_{n+k}}.
\end{equation*}
Note that in the limit, this gives (subject to convergence conditions)
\begin{equation} \label{Baileylimit}
	\lim_{n \to \infty} \beta_n = \frac{1}{(q,xq)_{\infty}} \sum_{k=0}^{\infty} \alpha_k.
\end{equation}
In practice, the above sum often becomes an infinite product by an appeal to the Jacobi triple product identity \cite[equation~(2.2.10)]{AndBook},
\begin{align}
	\sum_{n\in\ZZ}(-1)^n z^nq^{n^2}=\bigl(q^2,zq,z^{-1}q;q^2\bigr)_\infty,
	\label{eqn:jtp}
\end{align}
or the quintuple product identity \cite{C-quint},
\begin{align}
	\Biggl(\sum_{\substack{n\in\ZZ\\ n\equiv 0\,\,(\mathrm{mod}\,\,3)}}
	-\sum_{\substack{n\in\ZZ\\ n\equiv 2\,\,(\mathrm{mod}\,\,3)}}
	\Biggr) z^nq^{\frac{1}{3}\binom{n+1}{2}}=\bigl(q,zq,z^{-1}\bigr)_\infty\bigl(z^2q,z^{-2}q;q^2\bigr)_\infty
	\label{eqn:qtp}.
\end{align}

The most important fact about Bailey pairs is the following, which is known as the Bailey lemma. From a given Bailey pair relative to $x$, it produces new Bailey pairs relative to $x$.

\begin{Lemma}[\cite{Andmultiple}] \label{Baileylemma}
	If $(\alpha_n,\beta_n)$ is a Bailey pair relative to $x$, then so is $(\alpha_n',\beta_n')$, where
	\begin{equation*} %\label{Baileylemmaalpha}
		\alpha_n' = \frac{(b,c)_n(xq/bc)^n}{(xq/b,xq/c)_n}\alpha_n
	\qquad
	\text{and}
	\qquad %\label{Baileylemmabeta}
		\beta_n' = \frac{1}{(xq/b,xq/c)_n} \sum_{j=0}^n \frac{(b,c)_j(xq/bc)_{n-j} (xq/bc)^j}{(q)_{n-j}} \beta_j.
	\end{equation*}
\end{Lemma}

A result similar to Lemma \ref{Baileylemma} takes a Bailey pair relative to $x$ and produces Bailey pairs relative to $x/q$.
\begin{Lemma}[\cite{AABLattice}] \label{Baileylatticelemma}
	If $(\alpha_n,\beta_n)$ is a Bailey pair relative to $x$, then $(\alpha_n',\beta_n')$ is a Bailey pair relative to $x/q$, where
	\begin{equation*} %\label{Baileylatticealpha}
		\alpha_n' = (1-x) \Bigl(\frac{x}{bc} \Bigr) \frac{(b,c)_n(x/bc)^n}{(x/b,x/c)_n}\biggl(\frac{\alpha_n}{1-xq^{2n}} - \frac{xq^{2n-2}\alpha_{n-1}}{1-xq^{2n-2}} \biggr)
	\end{equation*}
and
	\begin{equation*} %\label{Baileylatticebeta}
		\beta_n' = \frac{1}{(x/b,x/c)_n} \sum_{j=0}^n \frac{(b,c)_j(x/bc)_{n-j} (x/bc)^j}{(q)_{n-j}} \beta_j.
	\end{equation*}
Here by convention we take
	\begin{equation} \label{eqn:latticeextracond}
		\alpha_{-1} = 0.
	\end{equation}
\end{Lemma}


Finally, we have two lemmas that change a Bailey pair relative to $x$ to one relative to $xq$. The first will be key to introducing $q$-binomial coefficients of the form $\bigl[\begin{smallmatrix} n_{i+1} + 1 \\ n_i \end{smallmatrix}\bigr]$. The second is not needed for the proof of Theorem \ref{CLSXYthm} but will be used to establish a certain Bailey pair later in the paper.
\begin{Lemma}[\cite{Lo1}] \label{littlelemma}
	If $(\alpha_n,\beta_n)$ is a Bailey pair relative to $x$, then $(\alpha_n',\beta_n')$ is a Bailey pair relative to $xq$, where
	\begin{equation*} %\label{littlelemmaalpha}
		\alpha_n' = \frac{1}{1-xq} \biggl(\frac{1-q^{n+1}}{1-xq^{2n+2}} \alpha_{n+1} + \frac{q^n(1-xq^n)}{1-xq^{2n}} \alpha_n \biggr)
	\qquad
	\text{and}
	\qquad %\label{littlelemmabeta}
		\beta_n' = \bigl(1-q^{n+1}\bigr)\beta_{n+1}.
	\end{equation*}
\end{Lemma}

\setcounter{section}{3}\setcounter{Theorem}{0}\setcounter{equation}{0}\subsection*{Author finite Bailey-pair transformation}
% Original source lines 326--459
\begin{Lemma} \label{lemma:mainpair}
	Let $m \geq 1$ and $0 \leq a \leq m$. Suppose that $(\alpha_n,\beta_n)$ is a Bailey pair relative to $q$ and that there is a sequence $(f_n)_{n \geq -1}$ such that
	\begin{align} \label{eqn:fdeflater}
		\alpha_n=\dfrac{1-q^{2n+1}}{1-q}f_n,
	\end{align}
	where
	\begin{align}
		q^{-1}f_{-1}=-f_0.
		\label{eqn:f-1f0}
	\end{align}
	Then $(\alpha_n',\beta_n')$ is a Bailey pair relative to $q$, where
	\begin{equation*} %\label{eqn:mainpairalpha}
		\alpha_n' =
		\begin{cases}
			\dfrac{1}{1-q}q^{m\binom{n+1}{2}}\bigl(
			q^{a(n+1)}f_{n+1} -q^{-an+2n-1}f_{n-1}
			\bigr) & \text{if $0 \leq a \leq m-1$}, \vspace{1mm}\\
			\dfrac{1}{1-q}q^{m \binom{n+1}{2}}\bigl(1-q^{2n+1}\bigr)f_n &\text{if $a=m$},
		\end{cases}
	\end{equation*}
	and
	\begin{align}
		\beta_n' &{}= \beta_{n_{m+1}}' \nonumber \\
 &{}= \frac{1}{(-q)_{n_{m+1}}} \sum_{n_m, \dots, n_1 \geq 0} \frac{q^{\binom{n_m+1}{2} + \cdots + \binom{n_1+1}{2}}\bigl(q^2;q^2\bigr)_{n_1+\delta_{a,0}}}{(q)_{n_{m+1}-n_m} (q)_{n_m}} \beta_{n_1+\delta_{a,0}} \prod_{i=1}^{m-1} \begin{bmatrix} n_{i+1} + \delta_{a,i} \\ n_i \end{bmatrix}.\label{eqn:mainpairbeta}
	\end{align}
	
\end{Lemma}

\begin{proof}
	{\it Step 1\emph{:}} First assume $a > 0$. Iterate Lemma \ref{Baileylemma} $a$ times with $x\mapsto q$, $b\mapsto -q$, and $c\rightarrow \infty$ to obtain another Bailey pair relative to $q$. The resulting $\alpha_n$ and $\beta_n$ are
	\begin{align}
		\alpha_n=q^{a\binom{n+1}{2}}\dfrac{1-q^{2n+1}}{1-q}f_n
		\label{eqn:alphafn}
	\end{align}
	and
	\begin{align}
		\beta_n=\beta_{n_{a+1}}=
		\dfrac{1}{(-q)_{n_{a+1}}}\sum_{n_a,\dots,n_1\geq 0}
		\dfrac{q^{\binom{n_a+1}{2}+\cdots+\binom{n_1+1}{2}}(-q)_{n_1}}{(q)_{n_{a+1}-n_a}\cdots (q)_{n_2-n_1}}\beta_{n_1}. \label{eqn:betastep1}
	\end{align}
	Here and throughout the paper we use the convention that $1/(q)_n = 0$ if $n <0$.
	If $a=m$, we stop here. Multiplying the numerator and denominator of the above sum by $(q)_{n_1} \cdots (q)_{n_m}$ and rewriting in terms of $q$-binomial coefficients using \eqref{eqn:qbinom} gives \eqref{eqn:mainpairbeta} for $a=m$.
	
	{\it Step 2\emph{:}} Apply Lemma \ref{littlelemma} with $x \mapsto q$ to obtain a Bailey pair relative to $q^2$:
	\begin{align*}
		\alpha_n &{}= \dfrac{1-q^{n+1}}{(q)_2}\bigl(
		q^{a\binom{n+2}{2}}f_{n+1} + q^{n+a\binom{n+1}{2}}f_n\bigr)\\
		&{}=\dfrac{1-q^{n+1}}{(q)_2}q^{(a+2)\binom{n+1}{2}-n^2}
		\bigl(
		f_n+q^{a(n+1)-n}f_{n+1}\bigr)
	\end{align*}
	and
	\begin{align*}
		\beta_n=\beta_{n_{a+1}}=
		\dfrac{1-q^{n_{a+1}+1}}{(-q)_{n_{a+1}+1}}\sum_{n_a,\dots,n_1\geq 0}
		\dfrac{q^{\binom{n_a+1}{2}+\cdots+\binom{n_1+1}{2}}(-q)_{n_1}}{(q)_{n_{a+1}+1-n_a}\cdots (q)_{n_2-n_1}}\beta_{n_1}.
	\end{align*}
	
	
	{\it Step 3\emph{:}} Use Lemma \ref{Baileylemma} with $x\mapsto q^2$, $b\mapsto-q^2$, and $c \to \infty$ to obtain another Bailey pair relative to $q^2$. The resulting pair is
	\begin{align*}
		\alpha_n
		=\dfrac{1-q^{2n+2}}{(1+q)(q)_2}q^{(a+3)\binom{n+1}{2}-n^2}
		\bigl(
		f_n+q^{a(n+1)-n}f_{n+1}\bigr)
	\end{align*}
	and
	\begin{align*}
		\beta_n=\beta_{n_{a+2}}=
		\dfrac{1}{(1+q)(-q)_{n_{a+2}}}\sum_{n_{a+1},\dots,n_1\geq 0}
		\dfrac{q^{\binom{n_{a+1}+1}{2}+\cdots+\binom{n_1+1}{2}}\bigl(1-q^{n_{a+1}+1}\bigr)(-q)_{n_1}}{(q)_{n_{a+2}-n_{a+1}}(q)_{n_{a+1}+1-n_a}\cdots (q)_{n_2-n_1}}\beta_{n_1}.
	\end{align*}
	At this point, we multiply both
	$\alpha_n$ and $\beta_n$ by $1+q$.
	
	
	{\it Step 4\emph{:}} Use Lemma \ref{Baileylatticelemma} with $x \mapsto q^2$ and $b,c \mapsto q$ to obtain a Bailey pair relative to $q$. This retains the $\beta_n$, but does change the $\alpha_n$. When $n \geq 1$, we have
	\begin{gather*}
		\alpha_n = \dfrac{1}{1-q}
		\bigl(
		q^{(a+3)\binom{n+1}{2}-n^2}\bigl(f_n+q^{a(n+1)-n}f_{n+1}\bigr) \\ \hphantom{\alpha_n = \dfrac{1}{1-q}\bigl(}{}
		- q^{2n+(a+3)\binom{n}{2}-(n-1)^2}\bigl(f_{n-1}+q^{an-n+1}f_{n}\bigr)
		\bigr)\\ \hphantom{\alpha_n}{}
		=
		\dfrac{1}{1-q}
		\bigl(
		q^{(a+3)\binom{n+1}{2}-n^2+a(n+1)-n}f_{n+1}
		-
		q^{2n+(a+3)\binom{n}{2}-(n-1)^2}f_{n-1}
		\bigr)\\ \hphantom{\alpha_n}{}
		= \dfrac{1}{1-q}q^{(a+1)\binom{n+1}{2}}\bigl(
		q^{a(n+1)}f_{n+1} -q^{-an+2n-1}f_{n-1}
		\bigr).
	\end{gather*}
	When $n=0$, the above is also valid, using \eqref{eqn:latticeextracond} and \eqref{eqn:f-1f0}.
	
	
	{\it Step 5\emph{:}} Iterate Lemma \ref{Baileylemma} $m-a-1$ times with $x\mapsto q$, $b\mapsto -q$, and $c\rightarrow \infty$ to arrive at the final Bailey pair relative to $q$:
	\begin{align*}
		\alpha_n
		=
		\dfrac{1}{1-q}q^{m\binom{n+1}{2}}\bigl(
		q^{a(n+1)}f_{n+1} -q^{-an+2n-1}f_{n-1}
		\bigr)
	\end{align*}
	and
	\begin{align*}
		\beta_n=\beta_{n_{m+1}}=
		\dfrac{1}{(-q)_{n_{m+1}}}\sum_{n_m,\dots,n_1\geq 0}
		\dfrac{q^{\binom{n_{m}+1}{2}+\cdots+\binom{n_1+1}{2}}\bigl(1-q^{n_{a+1}+1}\bigr)(-q)_{n_1}}
		{(q)_{n_{m+1}-n_{m}}\cdots
			(q)_{n_{a+1}+1-n_a}\cdots (q)_{n_2-n_1}}\beta_{n_1}.
	\end{align*}
	Rewriting the sum in terms of $q$-binomial coefficients gives{\samepage
	\begin{equation*} %\label{eqn:genbetafin}
		\beta_n = \beta_{n_{m+1}} = \frac{1}{(-q)_{n_{m+1}}} \sum_{n_m, \dots, n_1 \geq 0} \frac{q^{\binom{n_m+1}{2} + \cdots + \binom{n_1+1}{2}}\bigl(q^2;q^2\bigr)_{n_1}}{(q)_{n_{m+1}-n_m} (q)_{n_m}} \beta_{n_1} \prod_{i=1}^{m-1} \begin{bmatrix} n_{i+1} + \delta_{a,i} \\ n_i \end{bmatrix},
	\end{equation*}
	which coincides with \eqref{eqn:mainpairbeta}.}
	
	{\it Step 6\emph{:}}
	If $a=0$, performing Steps 2--5 leads to the following $\beta_n$ (while the formula for $\alpha_n$ is the same as in Step 5 with $a\mapsto 0$):
	\begin{align*}
		\beta_n=\beta_{n_{m+1}}=
		\dfrac{1}{(-q)_{n_{m+1}}}\sum_{n_m,\dots,n_1\geq 0}
		\dfrac{q^{\binom{n_{m}+1}{2}+\cdots+\binom{n_1+1}{2}}\bigl(1-q^{n_{1}+1}\bigr)(-q)_{n_1+1}}
		{(q)_{n_{m+1}-n_{m}}\cdots (q)_{n_2-n_1}}
		\beta_{n_1+1}.
	\end{align*}
	Rewriting this in terms of $q$-binomial coefficients gives
	\begin{equation*} %\label{eqn:genbetafina0}
		\beta_n = \beta_{n_{m+1}} = \frac{1}{(-q)_{n_{m+1}}} \sum_{n_m, \dots, n_1 \geq 0} \frac{q^{\binom{n_m+1}{2} + \cdots + \binom{n_1+1}{2}}\bigl(q^2;q^2\bigr)_{n_1+1}}{(q)_{n_{m+1}-n_m} (q)_{n_m}} \beta_{n_1+1} \prod_{i=1}^{m-1} \begin{bmatrix} n_{i+1} + \delta_{a,i} \\ n_i \end{bmatrix},
	\end{equation*}
	which again matches \eqref{eqn:mainpairbeta}. This completes the proof.
\end{proof}
\begin{thebibliography}{99}
\bibitem[1]{AABLattice}
Agarwal A.K., Andrews G.E., Bressoud D.M., The {B}ailey lattice,
 \textit{J.~Indian Math. Soc. (N.S.)} \textbf{51} (1987), 57--73.

\bibitem[3]{Andmultiple}
Andrews G.E., Multiple series {R}ogers--{R}amanujan type identities,
 \href{https://doi.org/10.2140/pjm.1984.114.267}{\textit{Pacific~J.~Math.}} \textbf{114} (1984), 267--283.

\bibitem[4]{AndBook}
Andrews G.E., The theory of partitions, \textit{Encyclopedia Math. Appl.}, \href{https://doi.org/10.1017/CBO9780511608650}{Cambridge
 University Press}, Cambridge, 1994.

\bibitem[9]{C-quint}
Cooper S., The quintuple product identity, \href{https://doi.org/10.1142/S1793042106000401}{\textit{Int.~J.~Number Theory}}
 \textbf{2} (2006), 115--161.

\bibitem[16]{Lo1}
Lovejoy J., Bailey pairs and strange identities, \href{https://doi.org/10.4134/JKMS.j220167}{\textit{J.~Korean Math. Soc.}}
 \textbf{59} (2022), 1015--1045, \href{https://arxiv.org/abs/2203.14391}{arXiv:2203.14391}.

\end{thebibliography}
\end{document}
